\section{Related Works}

\paragraph{Speculative decoding.} Speculative decoding can improve LLM inference latency without compromising the quality of the generated text. Other model-based methods include Medusa~\citep{cai2024medusasimplellminference} and SpecInfer~\citep{specinfer}, which first introduced tree-based speculation. Medusa uses multiple decoding heads for parallel candidate generation, while SpecInfer uses multiple draft models or speculative sampling to generate a tree of candidate tokens.
There are also several other model-free speculation methods like
LLMA~\citep{yang2023inferencereferencelosslessacceleration} and ANPD~\citep{ou2024losslessaccelerationlargelanguage}. These methods rely on small reference texts such as the prompt or retrieved contexts, and lack the adaptive speculation mechanisms of SuffixDecoding. Compared to prior speculative decoding methods, SuffixDecoding is uniquely designed for the emerging class of agentic LLM applications.

\paragraph{Methods for accelerating LLM agents.} Recent works also target the problem of latency in agentic applications. ALTO~\citep{10.1145/3642970.3655844} takes a systems-oriented approach to optimize the latency of multi-agent workflows through more efficient pipelining and scheduling. Dynasor~\citep{fu2024efficientlyservingllmreasoning} monitors the certainty of agentic reasoning algorithms as they run, and early-terminates reasoning paths that are unlikely to improve the final answer. Compared to these systems, SuffixDecoding takes an orthogonal speculative decoding approach, and they can be used in combination.
